\documentclass[12pt]{article}
\usepackage[utf8]{inputenc}
\usepackage[T1]{fontenc}
\usepackage[margin=1in]{geometry}
\usepackage{setspace}
\usepackage{xcolor}
\usepackage[colorlinks=true, urlcolor=blue!60!black, linkcolor=blue!60!black]{hyperref}
\setlength{\parindent}{2em}
\setlength{\parskip}{1.25em}
\renewcommand{\baselinestretch}{1.0}
\begin{document}
\begin{center}
{\LARGE \textit{Egbert v. Boule}, 596 U.S. 482 (2022)\par}
\vspace{0.8em}
{\large Case Brief\par}
\end{center}
\vspace{1em}

\section{Memory Jogger}
An innkeeper roughed up by a Border Patrol agent in his own driveway could not sue under \textit{Bivens}: the Court rejected his Fourth Amendment claim 6--3 and his First Amendment retaliation claim 9--0. The five-Justice majority explained that a court may not extend \textit{Bivens} if there is any reason to think Congress is better equipped to create a remedy (\href{https://supreme.justia.com/cases/federal/us/596/21-147/}{Justia}).

\section{Detailed Case Facts}
Respondent Robert Boule owned and operated the Smuggler's Inn, a bed-and-breakfast in Blaine, Washington, whose backyard abutted the U.S.--Canada border. Boule drove a vehicle with a ``SMUGLER'' vanity license plate, occasionally drove guests across or near the border, and also served as a paid confidential informant for federal agents, including petitioner U.S. Border Patrol Agent Erik Egbert.

On March 20, 2014, Boule informed Agent Egbert that a Turkish national who had flown into Seattle from New York had scheduled transportation to Smuggler's Inn later that day. Suspecting an unlawful border crossing, Egbert waited near the inn and followed Boule's SUV into the inn's driveway when the guest arrived. Boule asked Egbert to leave his property. According to Boule, Egbert refused, became violent, lifted Boule off the ground, and threw him against the SUV and then to the ground, injuring his back and shoulder. After verifying that the Turkish guest was lawfully in the United States, Egbert and responding officers departed (the guest unlawfully entered Canada later that night).

Boule lodged an administrative grievance with Egbert's Customs and Border Protection (CBP) supervisors and filed an administrative claim under the Federal Tort Claims Act (FTCA). Boule alleged that Egbert retaliated by reporting Boule's ``SMUGLER'' license plate to the Washington State Department of Licensing and contacting the Internal Revenue Service (IRS) and other agencies to prompt an audit of Boule's business. Although the IRS audit cleared Boule of wrongdoing, he spent over \$5,000 in accounting fees. CBP ultimately denied Boule's FTCA claim and took no formal disciplinary action against Egbert.

\section{Procedural History}
In January 2017, Boule sued Agent Egbert in his individual capacity in the U.S. District Court for the Western District of Washington under \textit{Bivens v. Six Unknown Fed. Narcotics Agents}, 403 U.S. 388 (1971), asserting Fourth Amendment excessive-force and First Amendment retaliation claims. The District Court granted summary judgment for Egbert on both claims, refusing to extend \textit{Bivens}. The Ninth Circuit reversed (998 F.3d 370 (2021)), holding that both claims could proceed under \textit{Bivens}, and denied rehearing en banc over the dissent of twelve judges. The Supreme Court granted certiorari on November 5, 2021 (declining to grant review on Egbert's third question asking the Court to overrule \textit{Bivens} outright).

\section{Judicial Votes}
The Supreme Court \textbf{reversed} the Ninth Circuit, by a \textbf{6--3} vote on the Fourth Amendment excessive-force claim and \textbf{9--0} on the First Amendment retaliation claim. The case was argued March 2, 2022, and decided June 8, 2022 (\href{https://supreme.justia.com/cases/federal/us/596/21-147/}{Justia}; \href{https://www.law.cornell.edu/supremecourt/text/21-147}{Cornell LII}):

\begin{itemize}
    \item \textbf{Majority Opinion:} Justice Thomas delivered the opinion of the Court, joined by Chief Justice Roberts and Justices Alito, Kavanaugh, and Barrett.
    \item \textbf{Concurring Opinions:} Justice Gorsuch filed an opinion concurring in the judgment. Justice Sotomayor, joined by Justices Breyer and Kagan, concurred in the judgment as to the First Amendment claim.
    \item \textbf{Dissenting Opinion:} Justice Sotomayor, joined by Justices Breyer and Kagan, dissented as to the Fourth Amendment claim.
\end{itemize}

\section{Holding}
\textbf{Questions Presented:}

\begin{enumerate}
  \item Does an implied damages cause of action under \textit{Bivens} extend to a Fourth Amendment excessive-force claim against a U.S. Border Patrol agent performing immigration-enforcement functions near the international border?
  \item Does a \textit{Bivens} cause of action extend to a First Amendment retaliation claim against a federal officer?
\end{enumerate}

\textbf{Holding:} No to both questions. Reversing the Ninth Circuit, the Supreme Court held that \textit{Bivens} does not authorize an implied damages remedy for either Boule's Fourth Amendment excessive-force claim (6--3) or his First Amendment retaliation claim (9--0). Collapsing the traditional two-step \textit{Bivens} framework into a single separation-of-powers inquiry, the Court held that a federal court may not recognize a \textit{Bivens} cause of action whenever there is \textit{any} reason to think that Congress might be better equipped than the Judiciary to create a damages remedy.

\section{Analysis of Opinions}
\subsection{Majority Opinion (Justice Thomas)}

Writing for the Court, Justice Thomas emphasized that creating a federal damages cause of action is a quintessentially legislative function; in the 42 years since \textit{Bivens} (1971), \textit{Davis v. Passman}, 442 U.S. 228 (1979), and \textit{Carlson v. Green}, 446 U.S. 14 (1980), the Court had declined eleven consecutive times to imply a new constitutional tort. While \textit{Ziglar v. Abbasi}, 582 U.S. 120 (2017), and \textit{Hern\'andez v. Mesa}, 589 U.S. 93 (2020), articulated a two-step test---(1) whether the claim arises in a ``new context'' and (2) whether ``special factors'' counsel hesitation---Justice Thomas explained that those two steps ``often resolve to a single question: whether there is any reason to think that Congress might be better equipped to create a damages remedy.'' If there is ``even a single reason to pause,'' a court may not imply a remedy.

\textbf{1. Fourth Amendment Excessive-Force Claim:} Accepting the Ninth Circuit's concession that the claim arose in a new context, the Court identified two independent reasons foreclosing \textit{Bivens} relief:
\begin{itemize}
  \item \textbf{Border Security and National Security:} Special-factors analysis must be conducted at a systemic level---asking whether courts are competent to authorize damages suits against Border Patrol agents \textit{generally}, not just against Agent Egbert on these specific facts. Echoing \textit{Hern\'andez}, the Court rejected the Ninth Circuit's granular distinctions (that Egbert stood on U.S. soil, investigated a guest who arrived from New York, and operated several feet from Canada), concluding that regulating Border Patrol conduct inevitably implicates national and cross-border security.
  \item \textbf{Alternative Remedial Structures:} Independently, Congress and the Executive had already provided an administrative grievance process within CBP (8 U.S.C. \S~1103(a)(3); 8 C.F.R. \S~287.10). Because \textit{Bivens} is concerned solely with deterring unconstitutional officer conduct rather than compensating victims, an administrative grievance mechanism forecloses \textit{Bivens} even if it affords the complainant no right to participate, no judicial review, and no money damages.
\end{itemize}

\textbf{2. First Amendment Retaliation Claim:} The Court unanimously refused to extend \textit{Bivens} to First Amendment retaliation claims. Because a plaintiff can easily allege retaliatory motive in response to any adverse federal action, authorizing damages suits would impose substantial social and litigation costs across the entire federal workforce, chilling legitimate official conduct. Furthermore, the Court rejected Boule's analogy to \textit{Davis v. Passman}: ``parallel circumstances'' with \textit{Bivens}, \textit{Passman}, or \textit{Carlson} are insufficient unless the claim also satisfies the restrictive analytical framework of the last four decades. Finally, although noting that the Court would not imply \textit{any} constitutional damages action if writing on a clean slate, the majority declined to overrule \textit{Bivens} itself.

\subsection{Concurring in the Judgment (Justice Gorsuch)}

Justice Gorsuch argued that the Court's case-by-case ``special factors'' inquiry leaves litigants with false hope because, under separation-of-powers principles, \textit{weighing the costs and benefits of new causes of action} is always Congress's job, never the Judiciary's. Seeing no principled difference between Boule's Fourth Amendment claim and \textit{Bivens} itself, he urged the Court to ``forthrightly'' overrule \textit{Bivens} and return the creation of federal damages remedies exclusively to Congress.

\subsection{Concurring in the Judgment in Part and Dissenting in Part (Justice Sotomayor)}

Justice Sotomayor agreed that Boule's First Amendment retaliation claim failed under \textit{Ziglar}'s two-step test because the Court had never recognized a First Amendment \textit{Bivens} action and retaliatory-motive suits risk chilling official action across every sphere of government. She sharply dissented, however, as to the Fourth Amendment claim, calling it ``materially indistinguishable'' from \textit{Bivens}: a federal law-enforcement officer entered a U.S. citizen's property on U.S. soil without a warrant and used excessive force. She accused the ``restless and newly constituted Court'' of rewriting \textit{Ziglar}, stretching national-security deference to immunize Border Patrol agents operating up to 100 miles from any U.S. border (where nearly 200 million Americans live), and treating a toothless internal CBP grievance process---which yielded ``no remedy whatsoever''---as an adequate alternative.

\subsection{Subsequent Developments}

\begin{itemize}

\item \textbf{\textit{Goldey v. Fields}, 606 U.S. \_\_\_ (June 30, 2025) (per curiam):} In a summary reversal of the Fourth Circuit (\href{https://www.supremecourt.gov/opinions/24pdf/24-809_9o6b.pdf}{SupremeCourt.gov}; \href{https://www.law.cornell.edu/supremecourt/text/24-809}{Cornell LII}; \href{https://supreme.justia.com/cases/federal/us/606/24-809/}{Justia}), the Supreme Court unanimously applied \textit{Egbert} to hold that no \textit{Bivens} remedy exists for a federal prisoner's Eighth Amendment excessive-force claim against Bureau of Prisons (BOP) officers. Quoting \textit{Egbert}, the Court underscored that Congress has legislated extensively in prisoner litigation (the Prison Litigation Reform Act) without creating a damages remedy, that systemic prison-administration concerns counsel hesitation, and that BOP's internal administrative grievance structure forecloses \textit{Bivens} even if ``not as effective as an individual damages remedy.''

\end{itemize}

\section{Examples: Future Applications}
\begin{enumerate}

\item \textbf{Same-Side (No \textit{Bivens} Remedy): An FBI Raid Far From the Border.} Suppose FBI agents execute a warrantless raid on a U.S. citizen's apartment in Houston, Texas (more than 100 miles from an international land border), handcuff the homeowner in front of his family, and break his arm without probable cause. He sues under \textit{Bivens}. The plaintiff will argue his claim falls squarely within \textit{Bivens}'s core domestic Fourth Amendment facts and involves no border-security factor. Under \textit{Egbert} and \textit{Goldey}, however, the agents will argue that FBI agents are a ``new category of defendants,'' that DOJ's Office of the Inspector General provides an ``alternative remedial structure,'' and that there is at least ``a single reason'' Congress is better equipped to authorize damages---likely defeating the claim unless the court finds no meaningful difference from \textit{Bivens}.

\item \textbf{Same-Side (No \textit{Bivens} Remedy): Retaliation by a Federal Inspector.} A restaurant owner publicly criticizes a federal food-safety agency. Days later, an agency inspector conducts an unannounced inspection and issues a string of citations that are later dismissed. The owner sues the inspector for damages under \textit{Bivens}, alleging First Amendment retaliation. \textit{Egbert} squarely forecloses the claim: all nine Justices agreed that Boule's retaliation claim failed. The majority declined to recognize a \textit{Bivens} action for First Amendment retaliation because a plaintiff can easily allege a retaliatory motive for almost any adverse federal action, and permitting such suits would chill legitimate official conduct across the federal workforce. Whether to provide a damages remedy for retaliation by federal officers is for Congress (\href{https://supreme.justia.com/cases/federal/us/596/21-147/}{Justia}).

\item \textbf{Opposite-Side (Express Statutory Cause of Action): The Same Force by a State Trooper.} A state highway patrol trooper stops a driver, drags him from the car without provocation, and throws him to the pavement, breaking his collarbone. The driver sues the trooper for damages under 42 U.S.C.\ \S~1983. \textit{Egbert} poses no obstacle. Its reasoning rests on separation of powers: courts should not \emph{imply} a damages action that Congress has not created. Here Congress has expressly supplied the cause of action against persons acting under color of state law, so there is no question whether Congress or the courts should decide. The driver may proceed on his Fourth Amendment excessive-force claim, subject to ordinary defenses such as qualified immunity.

\item \textbf{Opposite-Side (Prospective Relief): Enjoining an Ongoing Federal Practice.} Residents of a border town show that Border Patrol agents, under a written sector policy, routinely enter fenced private backyards without warrants or consent to look for migrants. The residents sue the agency head in her official capacity, seeking a declaratory judgment and an injunction against the practice; they seek no damages. \textit{Egbert} limits only implied retrospective money damages against individual officers. It does not withdraw the federal courts' power to grant prospective equitable relief against ongoing unconstitutional federal conduct, whether through traditional equity jurisdiction or under the Administrative Procedure Act's waiver of sovereign immunity for nonmonetary relief (5 U.S.C.\ \S~702). The suit may proceed to the merits.

\item \textbf{Fence-Sitter (A Near-Replica of \textit{Bivens}): Narcotics Agents Search a Home.} Federal narcotics agents assigned to a joint federal--state task force arrest a tenant in his New York apartment without a warrant or probable cause, manacle him in front of his family, search the apartment, and strip-search him at a federal facility. The tenant sues the federal agents for damages under the Fourth Amendment. The alleged search and arrest closely track \textit{Bivens}, and the plaintiff can argue that the agents acted in their ordinary federal narcotics-enforcement role. If the task force gave the agents a materially different legal mandate or supplied a distinct remedial process, a court could find a new context and decline to extend \textit{Bivens}. Whether the task-force assignment made a meaningful difference would turn on the facts.

\end{enumerate}

\section{Critique}
\textbf{Remedial Asymmetry (State vs. Federal Officers):} Under 42 U.S.C. \S~1983, state and local police officers face damages liability for Fourth Amendment excessive force and First Amendment retaliation. After \textit{Egbert}, federal officers committing identical constitutional violations---even on joint federal-state task forces---are largely shielded from personal constitutional damages unless Congress acts.

\textbf{Executive Self-Immunization:} Because \textit{Egbert} treats an agency's own internal complaint process as an ``alternative remedial structure'' regardless of whether it offers compensation or judicial review, executive agencies can arguably immunize their officers from \textit{Bivens} simply by maintaining a comment box or inspector-general hotline.

\textbf{Candor vs. Stealth Overruling:} Justice Gorsuch's concurrence exposes a sharp methodological tension: if the majority's single-question test can almost never be satisfied, retaining \textit{Bivens} on the books invites costly lower-court litigation over hair-splitting ``new context'' distinctions rather than cleanly overruling \textit{Bivens} or reaffirming its core.

\section{Key Quotations}
\begin{quote}
``Because our cases have made clear that, in all but the most unusual circumstances, prescribing a cause of action is a job for Congress, not the courts, we reverse.''
\end{quote}
\noindent --- Justice Thomas, for the Court, 596 U.S. at 486 (\href{https://www.law.cornell.edu/supremecourt/text/21-147}{Cornell LII}; \href{https://www.supremecourt.gov/opinions/21pdf/596us1r37_7l48.pdf}{SupremeCourt.gov}).

\begin{quote}
``While our cases describe two steps, those steps often resolve to a single question: whether there is any reason to think that Congress might be better equipped to create a damages remedy.''
\end{quote}
\noindent --- Justice Thomas, for the Court, 596 U.S. at 492 (\href{https://www.law.cornell.edu/supremecourt/text/21-147}{Cornell LII}; \href{https://supreme.justia.com/cases/federal/us/596/21-147/}{Justia}).

\begin{quote}
``So long as Congress or the Executive has created a remedial process that it finds sufficient to secure an adequate level of deterrence, the courts cannot second-guess that calibration by superimposing a \textit{Bivens} remedy.''
\end{quote}
\noindent --- Justice Thomas, for the Court, 596 U.S. at 498 (\href{https://www.law.cornell.edu/supremecourt/text/21-147}{Cornell LII}; \href{https://supreme.justia.com/cases/federal/us/596/21-147/}{Justia}).

\begin{quote}
``The Court goes to extraordinary lengths to avoid this result: It rewrites a legal standard it established just five years ago, stretches national-security concerns beyond recognition, and discerns an alternative remedial structure where none exists.''
\end{quote}
\noindent --- Justice Sotomayor, concurring in the judgment in part and dissenting in part, 596 U.S. at 505 (\href{https://www.supremecourt.gov/opinions/21pdf/596us1r37_7l48.pdf}{SupremeCourt.gov}; \href{https://www.law.cornell.edu/supremecourt/text/21-147}{Cornell LII}).


\end{document}
